%% Plot styles for the bisimulation figures.
%% The data is produced by scripts/postprocess.py (plotdata.csv).

%% Colours: the bright qualitative scheme of Paul Tol (T-Q-B) from the colorblind package.
\colorlet{tolblue}{T-Q-B1}
\colorlet{tolcyan}{T-Q-B2}
\colorlet{tolgreen}{T-Q-B3}
\colorlet{tolyellow}{T-Q-B4}
\colorlet{tolred}{T-Q-B5}
\colorlet{tolpurple}{T-Q-B6}
\colorlet{tolgrey}{T-Q-B0}

%% The colours of the configurations, which are the same in all figures. They differ from the
%% colours of the kinds of properties (blue and red).
\colorlet{colbisim}{tolpurple}   % exact bisimulation
\colorlet{colepsa}{tolcyan}      % tolerance 0.05
\colorlet{colepsb}{tolyellow}    % tolerance 0.1
\colorlet{colepsc}{tolgreen}     % tolerance 0.2

%% The IMDP that the plots are about, as it occurs in the csv columns (e.g. imdpa for the IMDPs learned
%% with identifier a, see mdp_to_imdp.py). Redefine it to switch to another IMDP.
\newcommand{\imdp}{imdpa}

%% Where the csv files written by postprocess.py live.
\newcommand{\resultsdir}{results}
\newcommand{\plotdata}{\resultsdir/plotdata.csv}

%% The models whose number of states the axes show: the input model, the MDP quotient and the IMDP
%% quotient, where the optional argument is the tolerance of approximate bisimulation.
\newcommand{\statesinput}{|S|}
\newcommand{\statesmdp}{\mathrm{MDP}/{\sim}}
\newcommand{\statesimdp}[1][]{\mathrm{IMDP}/{\sim}\ifthenelse{\equal{#1}{}}{}{_{\varepsilon=#1}}}

%% Smallest value in the data for the quantile plots; this must match PLOT_MIN_VALUE in postprocess.py.
\newcommand{\plotminvalue}{1e-8}
%% Position of the value 0 in the scatter plot of the values, which has logarithmic axes.
\newcommand{\plotzero}{1e-12}

%% Position of runs that exceeded the time or memory limit in the plots of the runtimes;
%% this must match PLOT_TIMEOUT in postprocess.py. \plotmaxtime is the upper end of the time axes.
\newcommand{\plottimeout}{20000}
\newcommand{\plotmaxtime}{45000}
%% Width of the quantile plot (without labels); chosen such that it fills a row together with a scatter plot.
\newlength{\quantileplotwidth}
\setlength{\quantileplotwidth}{0.545\linewidth}
\newlength{\scatterplotsize}
\setlength{\scatterplotsize}{0.222\linewidth}


%% The marks for the relation of an IMDP quotient to the MDP quotient (columns quotient-relation-*).
\tikzset{
  relation equal/.style={mark=*,tolgreen,mark size=1.3},
  relation finer/.style={mark=+,tolblue,mark size=1.9,semithick},
  relation coarser/.style={mark=x,tolpurple,mark size=2.1,thick},
  relation incomparable/.style={mark=star,tolred,mark size=1.9,semithick},
}
%% The marks for the kinds of properties: probabilities (P) and rewards (R), where the step-bounded
%% (finite horizon) variants have hollow marks.
\tikzset{
  property P indef/.style={mark=*,tolblue,mark size=1.3},
  property P fin/.style={mark=o,tolblue,mark size=1.4,semithick},
  property R indef/.style={mark=triangle*,tolred,mark size=1.7},
  property R fin/.style={mark=triangle,tolred,mark size=1.7,semithick},
  %% The marks for exact and approximate bisimulation in the plot of the quotienting times.
  %% The approximate one additionally gets the colour of its configuration.
  bisimulation exact/.style={mark=square,colbisim,mark size=1.3,semithick},
  bisimulation approx/.style={mark=x,mark size=1.9,semithick},
}
%% The marks for the classes of the largest absolute error of a run (columns error-class-*), which must
%% match ERROR_CLASSES in postprocess.py. Runs without any error value are in the class none.
\tikzset{
  error low/.style={mark=*,tolblue,mark size=1.3},
  error small/.style={mark=diamond*,tolyellow,mark size=1.8},
  error medium/.style={mark=triangle*,tolred,mark size=1.8},
  error high/.style={mark=star,black,mark size=2.0,semithick},
  error none/.style={mark=x,tolgrey,mark size=1.6,semithick},
}
\pgfplotsset{
  %% Only keep the rows where column #1 has the value #2.
  discard if not/.style 2 args={
    x filter/.append code={%
      \edef\tempa{\thisrow{#1}}\edef\tempb{#2}%
      \ifx\tempa\tempb\else\def\pgfmathresult{inf}\fi},
  },
  %% Common style of all (small) plots.
  small plot/.style={
    scale only axis,
    axis x line=bottom,
    axis y line=left,
    unbounded coords=discard, filter discard warning=false,
    % The x label has a fixed distance to the axis, independent of the height of the tick labels.
    xlabel style={font=\scriptsize, at={(axis description cs:0.5,0)}, anchor=north, yshift=-11pt},
    ylabel style={font=\scriptsize, yshift=-6pt},
    xticklabel style={font=\scriptsize},
    yticklabel style={font=\scriptsize},
    grid=major,
    xminorticks=false, yminorticks=false,
    axis on top=false,
    clip mode=individual,  % marks are drawn after (i.e. on top of) all lines
  },
  %% Legend with one column in the top left corner of a plot.
  legend inside/.style={
    legend columns=1,
    legend cell align={left},
    legend style={font=\tiny, inner sep=1pt, row sep=-2pt, at={(0.03,0.97)}, anchor=north west},
  },
  %% Common style of the quantile plots.
  quantile plot/.style={
    small plot,
    height=\scatterplotsize,
    ylabel style={yshift=2pt},
    legend cell align={left},
    legend style={font=\tiny, inner sep=1.5pt, at={(0.02,0.97)}, anchor=north west,
                  /tikz/every even column/.append style={column sep=4pt}},
    legend image code/.code={\draw[##1] plot coordinates {(0cm,0cm) (0.3cm,0cm)};},
    every axis plot/.append style={thick, no marks},
  },
  %% Legend above a plot; it is not part of the bounding box as it may be wider than the plot.
  legend above/.style={
    legend columns=4,
    legend style={font=\tiny, inner sep=1.5pt, /tikz/every even column/.append style={column sep=4pt},
                  at={(0.5,1.08)}, anchor=south, overlay},
  },
}

%% Counting how often each relation occurs in a column of the data.
\pgfplotstableread[col sep=comma]{\plotdata}\plotdatatable
\newcount\relationcountequal
\newcount\relationcountfiner
\newcount\relationcountcoarser
\newcount\relationcountincomparable
%% #1: column with the relation
\newcommand{\countrelations}[1]{%
  \global\relationcountequal=0 \global\relationcountfiner=0
  \global\relationcountcoarser=0 \global\relationcountincomparable=0
  \pgfplotstableforeachcolumnelement{#1}\of\plotdatatable\as\cell{%
    \ifcsname relationcount\cell\endcsname
      \global\advance\csname relationcount\cell\endcsname by 1
    \fi}%
}
%% Counts all rows with a value in column #1 as equal.
\newcommand{\countallequal}[1]{%
  \global\relationcountequal=0
  \pgfplotstableforeachcolumnelement{#1}\of\plotdatatable\as\cell{%
    \ifthenelse{\equal{\cell}{nan}}{}{\global\advance\relationcountequal by 1}}%
}
%% The marks of the relations as boxes. They are drawn here, i.e. outside of any axis, as a picture
%% within an axis would inherit the settings of that axis.
\newcommand{\definerelationmark}[1]{%
  \expandafter\newsavebox\csname relationmark#1\endcsname
  \expandafter\sbox\csname relationmark#1\endcsname{%
    \tikz[baseline=-0.55ex]{\draw[relation #1] plot[only marks] coordinates {(0,0)};}}}
\definerelationmark{equal}
\definerelationmark{finer}
\definerelationmark{coarser}
\definerelationmark{incomparable}
%% The mark of relation #1 followed by its count.
\newcommand{\relationcount}[1]{%
  \makebox[1.1em][c]{\expandafter\usebox\csname relationmark#1\endcsname}%
  \makebox[1em][r]{\the\csname relationcount#1\endcsname}}
%% A box with the counts of the given relations in a corner of the plot.
%% #1: comma separated list of relations  #2: north east or north west
\newcommand{\relationcounts}[2]{%
  \gdef\relationcountrows{}%
  \foreach \relation in {#1}{%
    \ifx\relationcountrows\empty
      \xdef\relationcountrows{\noexpand\relationcount{\relation}}%
    \else
      \xdef\relationcountrows{\unexpanded\expandafter{\relationcountrows}\noexpand\\\noexpand\relationcount{\relation}}%
    \fi}%
  \ifthenelse{\equal{#2}{north east}}{\def\countsx{1}}{\def\countsx{0}}%
  \edef\countsnode{\noexpand\node[anchor=#2, font=\noexpand\tiny, inner sep=1pt, outer sep=1.5pt, align=left,
      fill=white, fill opacity=0.85, text opacity=1, draw=black!40, draw opacity=1, line width=0.2pt]
    at (rel axis cs:\countsx,1)
    {\unexpanded\expandafter{\relationcountrows}};}%
  \countsnode
}

%% The points of a scatter plot, with marks according to the relation to the MDP quotient.
%% One plot per relation, so that the rarer relations are drawn on top of the frequent ones.
%% #1: table options that select x and y  #2: column with the relation  #3: show legend (true/false)
\newcommand{\relationscatter}[3]{%
  \foreach \relation in {equal,finer,coarser,incomparable}{%
    \edef\plotpoints{\noexpand\addplot[only marks, relation \relation, discard if not={#2}{\relation}, forget plot]
      table [col sep=comma, \unexpanded{#1}] {\noexpand\plotdata};}%
    \plotpoints
    % The legend does not depend on which relations occur in this plot.
    \ifthenelse{\equal{#3}{true}}{%
      \edef\plotlegend{\noexpand\addlegendimage{only marks, relation \relation}\noexpand\addlegendentry{\relation}}%
      \plotlegend}{}%
  }%
}

%% A scatter plot comparing two state counts.
%% #1: x column  #2: x label  #3: y column  #4: y label
%% #5: column with the relation to the MDP quotient that determines the marks (empty: all points are equal)
%% #6: show legend (true/false)
\newcommand{\statesscatterplot}[6]{%
  \begin{tikzpicture}
    \begin{axis}[
      small plot,
      width=\scatterplotsize,
      height=\scatterplotsize,
      axis equal image,
      xmin=1e2, xmax=1e8, ymin=1e2, ymax=1e8,
      xmode=log, ymode=log,
      xtick={1e2,1e4,1e6,1e8},
      ytick={1e2,1e4,1e6,1e8},
      xlabel={#2}, ylabel={#4},
      legend above,
    ]
      \addplot[no marks, forget plot] coordinates {(1e2,1e2) (1e8,1e8)};
      \addplot[no marks, densely dotted, forget plot] coordinates {(1e3,1e2) (1e8,1e7)};
      \addplot[no marks, densely dotted, forget plot] coordinates {(1e4,1e2) (1e8,1e6)};
      \ifthenelse{\equal{#5}{}}{%
        \addplot[only marks, relation equal]
          table [col sep=comma, x={#1}, y={#3}] {\plotdata};
        \countallequal{#3}
        \relationcounts{equal}{north west}
      }{%
        \relationscatter{x={#1}, y={#3}}{#5}{#6}
        \countrelations{#5}
        \relationcounts{equal,finer,coarser,incomparable}{north west}
      }
    \end{axis}
  \end{tikzpicture}%
}

%% A scatter plot of the size of an IMDP quotient relative to the size of the MDP quotient:
%% a point (x,y) means that the MDP quotient has x states and the IMDP quotient has x*y states.
%% #1: label of the IMDP configuration in the csv columns (e.g. imdp-exact or imdp-0.05)
%% #2: the number of states of the IMDP quotient (math)
%% #3: show legend (true/false)
\newcommand{\quotientratioplot}[3]{%
  \begin{tikzpicture}
    \begin{axis}[
      small plot,
      width=\scatterplotsize,
      height=\scatterplotsize,
      xmin=1e2, xmax=1e8, ymin=0.3, ymax=500,
      xmode=log, ymode=log,
      xtick={1e2,1e4,1e6,1e8},
      ytick={1,10,100},
      yticklabels={$1$,$10$,$100$},
      xlabel={$\statesmdp$}, ylabel={$#2$ ratio},
      ylabel style={yshift=2pt},  % the tick labels are wider here
      legend above,
    ]
      \addplot[no marks, forget plot] coordinates {(1e2,1) (1e8,1)};
      \relationscatter{x=quotient-states-mdp-exact,
                       y expr=\thisrow{quotient-states-#1}/\thisrow{quotient-states-mdp-exact}}{quotient-relation-#1}{#3}
      \countrelations{quotient-relation-#1}
      \relationcounts{equal,finer,coarser,incomparable}{north east}
    \end{axis}
  \end{tikzpicture}%
}

%% The points of a scatter plot that compares two configurations on the properties of both kinds, with
%% marks according to the kind of property (P, P<=, R, R<=).
%% #1: prefix of the csv columns (the columns are <prefix>-<kind>-<configuration>)
%% #2: label of the x configuration in the csv columns  #3: label of the y configuration
%% #4: smaller values are shown at this position  #5: show legend (true/false)
\newcommand{\propertyscatter}[5]{%
  \foreach \class/\kind/\name in {P/indef/$\mathrm{P}$, P/fin/$\mathrm{P}^{\le}$, R/indef/$\mathrm{R}$, R/fin/$\mathrm{R}^{\le}$}{%
    \edef\plotpoints{\noexpand\addplot[only marks, property \class\space\kind, discard if not={property-class}{\class}]
      table [col sep=comma,
             x expr={max(\noexpand\thisrow{#1-\kind-#2},#4)},
             y expr={max(\noexpand\thisrow{#1-\kind-#3},#4)}] {\noexpand\plotdata};}%
    \plotpoints
    \ifthenelse{\equal{#5}{true}}{%
      \edef\plotlegend{\noexpand\addlegendentry{\unexpanded\expandafter{\name}}}\plotlegend}{}%
  }%
}

%% The lines of a quantile plot.
%% #1: csv file, relative to \resultsdir  #2: comma separated list of column/colour/legend-entry triples
\newcommand{\quantilelines}[2]{%
  \foreach \col/\colour/\name in {#2}{%
    \edef\plotline{\noexpand\addplot[\colour]
      table [x=i, y=\col, col sep=comma] {\resultsdir/#1};
      \noexpand\addlegendentry{\unexpanded\expandafter{\name}}}%
    \plotline
  }%
}

\newcount\propertycounter
%% A quantile plot of the absolute errors: a point (x,y) on a line means that there are x properties
%% for which the absolute error of the configuration is at most y. The x axis ends at the number of properties.
%% #1: comma separated list of column/colour/legend-entry triples
\newcommand{\errorquantileplot}[1]{%
  % The number of properties is the number of values of exact bisimulation.
  \pgfplotstableread[col sep=comma]{\resultsdir/quantile-absdiff.csv}\quantiletable
  \global\propertycounter=0
  \edef\countproperties{\noexpand\pgfplotstableforeachcolumnelement{\imdp-exact}}%
  \countproperties\of\quantiletable\as\cell{%
    \ifthenelse{\equal{\cell}{nan}}{}{\global\advance\propertycounter by 1}}%
  \edef\propertycount{\the\propertycounter}%
  \begin{tikzpicture}
    \begin{axis}[
      quantile plot,
      ymode=log,  % cannot be part of a style
      width=\quantileplotwidth,
      xmin=0, xmax=\propertycount,
      ymin=\plotminvalue, ymax=0.1,
      xtick={0,20,40,60,80},
      extra x ticks={\propertycount},
      ytick={1e-8,1e-7,1e-6,1e-5,1e-4,1e-3,1e-2,1e-1},
      yticklabels={${\le}10^{-8}$,,$10^{-6}$,,$10^{-4}$,,$10^{-2}$,},
      xlabel={properties (\propertycount{} in total)}, ylabel={absolute error},
      legend columns=1,
      restrict y to domain*=-20:0,  % the natural logarithm of the error; larger errors leave the plot at the top
    ]
      \quantilelines{quantile-absdiff.csv}{#1}
    \end{axis}
  \end{tikzpicture}%
}

%% A scatter plot comparing the values of the properties (of both kinds) for two configurations.
%% The value 0 is shown at \plotzero.
%% #1: label of the x configuration in the csv columns  #2: x label
%% #3: label of the y configuration in the csv columns  #4: y label
\newcommand{\valuescatterplot}[4]{%
  \begin{tikzpicture}
    \begin{axis}[
      small plot,
      width=\scatterplotsize,
      height=\scatterplotsize,
      xmin=\plotzero, xmax=1e8, ymin=\plotzero, ymax=1e8,
      xmode=log, ymode=log,
      xtick={1e-12,1e-8,1e-4,1,1e4,1e8},
      xticklabels={$0$,,$10^{-4}$,,$10^{4}$,},
      ytick={1e-12,1e-8,1e-4,1,1e4,1e8},
      yticklabels={$0$,,$10^{-4}$,,$10^{4}$,},
      xlabel={#2}, ylabel={#4},
      legend inside,
    ]
      \addplot[no marks, forget plot] coordinates {(\plotzero,\plotzero) (1e8,1e8)};
      \propertyscatter{value}{#1}{#3}{\plotzero}{true}
    \end{axis}
  \end{tikzpicture}%
}

%% A scatter plot of the time for bisimulation quotienting against the size of the input model,
%% for exact bisimulation of the IMDP and for approximate bisimulation with the given tolerance.
%% #1: tolerance of approximate bisimulation (as in the csv columns)  #2: colour of that configuration
\newcommand{\bisimulationtimeplot}[2]{%
  \begin{tikzpicture}
    \begin{axis}[
      small plot,
      width=\scatterplotsize,
      height=\scatterplotsize,
      xmin=1e2, xmax=1e8, ymin=1e-3, ymax=1e3,
      xmode=log, ymode=log,
      xtick={1e2,1e4,1e6,1e8},
      ytick={1e-3,1e-2,1e-1,1,1e1,1e2,1e3},
      yticklabels={$10^{-3}$,,$0.1$,,$10$,,$10^{3}$},
      xlabel={$\statesinput$}, ylabel={quotienting time},
      legend inside,
    ]
      \addplot[only marks, bisimulation exact]
        table [col sep=comma, x=states, y=bisimulation-time-\imdp-exact] {\plotdata};
      \addlegendentry{$\sim$}
      \addplot[only marks, bisimulation approx, #2]
        table [col sep=comma, x=states, y=bisimulation-time-\imdp-#1] {\plotdata};
      \addlegendentry{$\varepsilon{=}#1$}
    \end{axis}
  \end{tikzpicture}%
}

%% A scatter plot of the wallclock times of the whole run: on the IMDP against with a quotient of it.
%% The marks tell how large the absolute error is; larger errors are drawn on top.
%% #1: label of the quotient configuration in the csv columns  #2: y label
\newcommand{\wallclocktimeplot}[2]{%
  \begin{tikzpicture}
    \begin{axis}[
      small plot,
      width=\scatterplotsize,
      height=\scatterplotsize,
      xmin=0.1, xmax=\plotmaxtime, ymin=0.1, ymax=\plotmaxtime,
      xmode=log, ymode=log,
      xtick={0.1,1,10,100,1000},
      xticklabels={$0.1$,,$10$,,$10^{3}$},
      ytick={0.1,1,10,100,1000},
      yticklabels={$0.1$,,$10$,,$10^{3}$},
      extra x ticks={\plottimeout}, extra x tick labels={TO},
      extra x tick style={grid=major},
      extra y ticks={\plottimeout}, extra y tick labels={TO},
      extra y tick style={grid=major},
      xlabel={IMDP}, ylabel={#2},
    ]
      \addplot[no marks, forget plot] coordinates {(0.1,0.1) (7200,7200)};
      \addplot[no marks, densely dotted, forget plot] coordinates {(1,0.1) (7200,720)};
      \addplot[no marks, densely dotted, forget plot] coordinates {(0.1,1) (720,7200)};
      \foreach \class in {none,low,small,medium,high}{%
        \edef\plotpoints{\noexpand\addplot[only marks, error \class, discard if not={error-class-#1}{\class}, forget plot]
          table [col sep=comma, x=wallclock-time-\imdp, y=wallclock-time-#1] {\noexpand\plotdata};}%
        \plotpoints
      }
    \end{axis}
  \end{tikzpicture}%
}

%% The legend for the marks of \wallclocktimeplot, as a box that can be placed next to the plots.
\newcommand{\errorclasslegend}{%
  \begin{tikzpicture}[font=\tiny]
    \matrix[draw, inner sep=2pt, row sep=1pt, column sep=3pt, nodes={inner sep=0pt, anchor=west}] {
      \node {largest absolute error}; \\
      \draw[error low] plot[only marks] coordinates {(0,0)}; \node at (0.25,0) {${\le}\,10^{-3}$}; \\
      \draw[error small] plot[only marks] coordinates {(0,0)}; \node at (0.25,0) {${>}\,10^{-3}$}; \\
      \draw[error medium] plot[only marks] coordinates {(0,0)}; \node at (0.25,0) {${>}\,10^{-2}$}; \\
      \draw[error high] plot[only marks] coordinates {(0,0)}; \node at (0.25,0) {${>}\,10^{-1}$}; \\
    };
  \end{tikzpicture}%
}
